\section{Impulse limit and Doppler width}\label{app:D.32}

A nucleus of mass $M_\kappa$ at position $\vect{X}$ with momentum $\vect{p}$ takes up the momentum $\vect{q}$ through $\op{O} = \ee^{\ii\vect{q}\cdot\vect{X}}$. We derive the moments of its self response for any potential, its impulse limit, the harmonic width, the free-recoil rate of \cref{eq:11.4-LS} and the absorption rate \cref{eq:11.4-ferm-nuc}. Throughout, $E_R = q^2/2M_\kappa$ and $p_{\hat q} = \hat{\vect{q}}\cdot\vect{p}$.

\subsection{The self correlation}

The self part of \cref{eq:11.3-exact} for one nucleus is $c_\kappa^2\Omega^{-1}\int\dd t\,\ee^{\ii\omega t}I(t)$ with $I(t) = \langle\ee^{-\ii\vect{q}\cdot\vect{X}(t)}\ee^{\ii\vect{q}\cdot\vect{X}}\rangle$. Write $\ee^{-\ii\vect{q}\cdot\vect{X}(t)} = \ee^{\ii Ht}\ee^{-\ii\vect{q}\cdot\vect{X}}\ee^{-\ii Ht}$ and move the kick through the second evolution operator,
\begin{equation*}
  \ee^{-\ii\vect{q}\cdot\vect{X}}\,\ee^{-\ii Ht}\,\ee^{\ii\vect{q}\cdot\vect{X}} = \ee^{-\ii H_{\vect{q}}t} , \qquad H_{\vect{q}} = \ee^{-\ii\vect{q}\cdot\vect{X}}H\,\ee^{\ii\vect{q}\cdot\vect{X}} .
\end{equation*}
The translation $\ee^{-\ii\vect{q}\cdot\vect{X}}\vect{p}\,\ee^{\ii\vect{q}\cdot\vect{X}} = \vect{p} + \vect{q}$ changes the kinetic energy $\vect{p}^2/2M_\kappa$ and nothing else, provided the potential energy depends on $\vect{X}$ but not on $\vect{p}$. All Coulomb interactions are of this kind. Then
\begin{equation}
  I(t) = \bigl\langle\ee^{\ii Ht}\,\ee^{-\ii(H + D)t}\bigr\rangle , \qquad D = H_{\vect{q}} - H = E_R + \frac{\vect{q}\cdot\vect{p}}{M_\kappa} ,
  \label{eq:D.32-I}
\end{equation}
exactly, for any potential and any stationary state. The operator $D$ is the energy that a kick $\vect{q}$ would give a free particle of momentum $\vect{p}$.

\subsection{Moments of the recoil spectrum}

The moments of the spectrum are the derivatives of $I$ at $t = 0$, $\int\dd\omega\,\omega^kS/\int\dd\omega\,S = \ii^k\,\dd^kI/\dd t^k|_0$ with the sign of \cref{eq:4.2-Sab}. The first derivative gives $\langle(H + D) - H\rangle = \langle D\rangle = E_R$, since $\langle\vect{p}\rangle = 0$; this is \cref{eq:5.6-sum}. The second derivative gives $\langle H^2 - 2H(H + D) + (H + D)^2\rangle = \langle D^2\rangle + \langle[D, H]\rangle$, and the average of a commutator with $H$ vanishes in a stationary state. Hence the variance
\begin{equation}
  \sigma_\omega^2 = \langle D^2\rangle - E_R^2 = \frac{q^2\langle p_{\hat q}^2\rangle}{M_\kappa^2} = 4E_R\,E^{\mathrm{kin}}_\kappa(\hat{\vect{q}}) , \qquad E^{\mathrm{kin}}_\kappa(\hat{\vect{q}}) = \frac{\langle p_{\hat q}^2\rangle}{2M_\kappa} ,
  \label{eq:D.32-var}
\end{equation}
holds at any temperature and for any potential local in $\vect{X}$. The width of the response measures the kinetic energy of the nucleus along $\hat{\vect{q}}$. At $T = 0$, with $|0\rangle$ the ground state, the third moment is $\langle0|(H - E_0 + D)^3|0\rangle$. Every term with $H - E_0$ at either end vanishes, which leaves $\langle D^3\rangle + \langle D(H - E_0)D\rangle$. The first is $E_R^3 + 3E_R\sigma_\omega^2$, because $\langle p_{\hat q}^3\rangle = 0$ on a site with inversion symmetry or by time reversal. The second is half the double commutator $\langle[D,[H, D]]\rangle$, and with $[V, p_{\hat q}] = \ii\,\partial_{\hat q}V$ this is $(q^2/2M_\kappa^2)\langle\partial_{\hat q}^2V\rangle$, the mean curvature of the potential along $\hat{\vect{q}}$. Subtracting the powers of the lower moments, the third cumulant is
\begin{equation}
  c_3 = \frac{q^2}{2M_\kappa^2}\,\bigl\langle\partial_{\hat q}^2V\bigr\rangle .
  \label{eq:D.32-c3}
\end{equation}
For a nucleus whose electrons follow it, at $n_B = 0$, $V$ is the Born--Oppenheimer surface of \cref{sec:5.1}. Taken with the full Coulomb Hamiltonian instead, $V$ contains the attraction of the electrons, whose curvature at the nucleus is the electron density there; that third moment includes the electronic excitations that the kick causes and is the subject of \cref{sec:12.5}.

\subsection{The moments at any temperature}

In thermal equilibrium the response is the average over the initial eigenstates $|i\rangle$ of $H$ with the Boltzmann weights $p_i$. From one eigenstate, the moments of the spectrum of $\omega = E_f - E_i$ are $\sum_f|\langle f|\ee^{\ii\vect{q}\cdot\vect{X}}|i\rangle|^2(E_f - E_i)^k = \langle i|(A_i + D)^k|i\rangle$, with $A_i = H - E_i$, by the same translation of $H$ as above. For a Hamiltonian invariant under time reversal the eigenfunctions can be chosen real, and then $\langle i|p_{\hat q}|i\rangle = \langle i|p_{\hat q}^3|i\rangle = 0$: in a real state the expectation value of $p_{\hat q} = -\ii\partial_{\hat q}$, or of its cube, is $\ii$ times a real integral, while it must be real. The mean is therefore $E_R$ from every $|i\rangle$, and the central moments are $\langle i|(A_i + \delta D)^k|i\rangle$ with $\delta D = (q/M_\kappa)p_{\hat q}$. Since $A_i|i\rangle = 0$, every term with $A_i$ at either end vanishes, exactly as for the ground state. The second moment is $\langle i|\delta D^2|i\rangle$. The third is $\langle i|\delta D^3|i\rangle + \langle i|\delta D\,A_i\,\delta D|i\rangle$; the first term vanishes, and the second is $\frac12\langle i|[\delta D,[H,\delta D]]|i\rangle = (q^2/2M_\kappa^2)\langle i|\partial_{\hat q}^2V|i\rangle$. The fourth is $\langle i|\delta D^4|i\rangle + \|A_i\,\delta D|i\rangle\|^2$ plus the mixed term $\langle i|\delta D\,A_i\,\delta D^2|i\rangle$ and its complex conjugate. Since $A_i\,\delta D|i\rangle = [H, \delta D]|i\rangle = \ii(q/M_\kappa)\,\partial_{\hat q}V|i\rangle$, the norm is $(q/M_\kappa)^2\langle i|(\partial_{\hat q}V)^2|i\rangle$; the mixed term is $\ii(q/M_\kappa)^3\langle i|p_{\hat q}(p_{\hat q}\partial_{\hat q}V + \partial_{\hat q}V\,p_{\hat q})|i\rangle$, $\ii$ times a real number in a real state, and cancels against its conjugate. Averaging over $i$ with the weights $p_i$ gives the moments \cref{eq:11.5-moments} with thermal averages in place of ground-state ones: all five hold at any temperature. For one harmonic oscillator of frequency $\omega_0$ the fourth cumulant that results is $E_R\omega_0^3\coth(\omega_0/2T)$, the even case of \cref{eq:D.32-ck} below.

\subsection{The impulse approximation}

If $D$ commuted with $H$, \cref{eq:D.32-I} would be $I(t) = \langle\ee^{-\ii Dt}\rangle$. For a free particle this is exact. For a bound one the commutator $[H, D] = \ii(q/M_\kappa)\,\partial_{\hat q}V$ is the force along $\hat{\vect{q}}$, and dropping it is the impulse approximation: the nucleus is treated as free during the kick. Writing the average over the momentum distribution $\mathfrak{f}_\kappa(\vect{p})$ of the nucleus, normalized to one,
\begin{equation}
  I_{\mathrm{IA}}(t) = \int\dd^3p\;\mathfrak{f}_\kappa(\vect{p})\,\ee^{-\ii(E_R + \vect{q}\cdot\vect{p}/M_\kappa)t} , \qquad
  \int\dd t\;\ee^{\ii\omega t}I_{\mathrm{IA}}(t) = 2\pi\int\dd^3p\;\mathfrak{f}_\kappa(\vect{p})\,\delta\Bigl(\omega - E_R - \frac{\vect{q}\cdot\vect{p}}{M_\kappa}\Bigr) .
  \label{eq:D.32-IA}
\end{equation}
Its mean and variance are those of \cref{eq:D.32-var}, exactly, and its third cumulant vanishes. The approximation therefore errs first in $c_3$. In the scaled variable $(\omega - E_R)/\sigma_\omega$ the omitted skewness is $c_3/\sigma_\omega^3$. With $c_3 \propto q^2$ and $\sigma_\omega \propto q$, it falls as $1/q$, while every cumulant that \cref{eq:D.32-IA} keeps scales with the width. The impulse approximation is the leading term of an expansion in $1/q$ at fixed $(\omega - E_R)/\sigma_\omega$. The time scale makes the same point: the integral over $t$ is cut off at $t \sim 1/\sigma_\omega$, and the force changes the momentum in that time by an amount small compared with its spread when $\sigma_\omega$ is large compared with the vibration frequencies.

\subsection{The harmonic crystal}

In a harmonic crystal $\mathfrak{f}_\kappa$ is Gaussian. From \cref{eq:5.6-C},
\begin{equation*}
  \langle p^\alpha p^\beta\rangle = M_\kappa^2\langle\dot u^\alpha\dot u^\beta\rangle = \frac{M_\kappa}{2N}\sum_{\vect{Q}\nu}\omega_{\vect{Q}\nu}\,e^\alpha_{\kappa\nu}(\vect{Q})\,e^\beta_{\kappa\nu}(\vect{Q})^*\coth\frac{\omega_{\vect{Q}\nu}}{2T} .
\end{equation*}
On a site of cubic symmetry the tensor is diagonal, and with the projected density of states $g_{\mathrm{ph},\kappa}/3$ of \cref{sec:11.3}, normalized to one,
\begin{equation}
  E^{\mathrm{kin}}_\kappa = \frac14\int\dd\omega\;\frac{g_{\mathrm{ph},\kappa}(\omega)}{3}\;\omega\coth\frac{\omega}{2T} \;\xrightarrow{T = 0}\; \frac{\langle\omega\rangle}{4} ,
  \label{eq:D.32-Ekin}
\end{equation}
half the zero-point energy of each mode, averaged. The impulse spectrum \cref{eq:D.32-IA} is then a Gaussian of mean $E_R$ and variance $E_R\langle\omega\coth(\omega/2T)\rangle$. The same Gaussian follows from the generating formula \cref{eq:5.6-gen} directly. Its exponent $\vect{q}\cdot\mathcal{C}_{\kappa\kappa}(t)\cdot\vect{q}$, expanded to second order in $t$, is $2W_\kappa - \ii E_Rt - \frac12E_R\langle\omega\coth(\omega/2T)\rangle t^2$; the constant cancels the Debye--Waller factor, and the time integral is Gaussian. This is the saddle-point evaluation of Campbell-Deem and collaborators~\cite[Eqs.~(37)--(39) and App.~B]{campbelldeem2022dark}. Every order of the expansion is available: the $k$-th derivative of the exponent at $t = 0$ gives the cumulant
\begin{equation}
  c_k = E_R\,\bigl\langle\omega^{k-1}\bigr\rangle\ (k\ \text{odd}) , \qquad c_k = E_R\,\Bigl\langle\omega^{k-1}\coth\frac{\omega}{2T}\Bigr\rangle\ (k\ \text{even}) ,
  \label{eq:D.32-ck}
\end{equation}
which at $T = 0$ is \cref{eq:8.5-cumulants}. The odd cumulants do not depend on temperature. For $k = 3$, \cref{eq:D.32-c3} agrees, since the on-site force constant of a harmonic crystal is $M_\kappa\langle\omega^2\rangle$ by completeness of the polarization vectors.

\subsection{Free recoil and the standard rate}

As $x_R = E_R/\langle\omega\rangle$ grows, $\sigma_\omega/E_R = x_R^{-1/2}$ at $T = 0$ and the Gaussian tends to $\delta(\omega - E_R)$~\cite[Eq.~(42)]{campbelldeem2022dark}. For a coupling to nucleons the vertex of each nucleus is $A_\kappa F_N(q)$, the elastic form factor \cref{eq:2.2-FN}; elastic, because the momenta of the impulse regime stay far below the energies that excite the nucleus internally. For one species the response is then $S_A = 2\pi\nd_NA^2F_N(q)^2\delta(\omega - q^2/2M)$, with $\nd_N = n_{\mathrm{at}}/\Omega$ the number density of nuclei. Insert it into \cref{eq:8.2-spectrum} with one row, $\mathfrak{D} = |\tilde V_N|^2$. The response does not depend on $\hat{\vect{q}}$, so the angle between $\vect{q}$ and $\vect{v}$ can be integrated first: $\int\dd^3q/(2\pi)^3$ of $\delta(\omega - \vect{q}\cdot\vect{v} + q^2/2\mchi)$ gives $\int q\,\dd q\,\Theta(v - \vmin)/(4\pi^2v)$, as in \cref{sec:1.3}, and the average over $f_\chi$ turns $\Theta/v$ into $\eta(\vmin)$ of \cref{eq:1.3-eta}. With $\rho_T = \nd_NM$,
\begin{equation*}
  \frac{\dd\mathcal{R}}{\dd\omega} = \frac{\rho_\chi}{\mchi\,\nd_NM}\int\frac{q\,\dd q}{4\pi^2}\,|\tilde V_N|^2\,2\pi\nd_NA^2F_N^2\,\delta\Bigl(\omega - \frac{q^2}{2M}\Bigr)\,\eta(\vmin)
  = \frac{\rho_\chi A^2|\tilde V_N|^2}{2\pi\mchi}\,F_N(q)^2\,\eta(\vmin) ,
\end{equation*}
at $q = (2M\omega)^{1/2}$, and \cref{eq:1.3-vmin} with $\Delta_\chi = 0$ gives $\vmin = q/2M + q/2\mchi = q/2\mu_{\chi N}$. A heavy mediator, $|\tilde V_N|^2 = \pi\sigma_p/\mu_{\chi p}^2$, with $\sigma_N = \sigma_pA^2\mu_{\chi N}^2/\mu_{\chi p}^2$ from \cref{eq:2.2-nucleus}, gives \cref{eq:11.4-LS}. The same steps applied to the impulse Gaussian give the rate with the width, which is what \cref{fig:11-limit-4a} evaluates.

\subsection{Absorption on a nucleus}

For a fermion absorbed at rest the momentum transfer is $q = \mchi - \omega$ (\cref{eq:8.1-omega-ferm}), and the vector rate of \cref{eq:10.9-ferm} is $(\sigma_F/2\pi)\bar S_{nn}(q,\omega)$, where the density now counts nucleons. The two operators differ in the spin trace. With the dark matter at rest, $p = (\mchi, \vect{0})$, a massless outgoing fermion of four-momentum $(q, -\vect{q})$ and the nucleon current $J = (\rho, (\omega/q)\rho\,\hat{\vect{q}})$ fixed by continuity, the spin-summed contraction is $\mchi^3\rho^2/q$ for the vector operator and $\mchi q\rho^2$ for the scalar density; their ratio is $q^2/\mchi^2$. Continuity holds for the nucleon current exactly, and the transverse current of a nucleus is of relative order $E^{\mathrm{kin}}_\kappa/M$ or $q^2/M^2$ and is dropped. For a free nucleus,
\begin{equation*}
  \varGamma_V = \sigma_F\,\nd_NA^2\int\dd\omega\;F_N(\mchi - \omega)^2\,\delta\Bigl(\omega - \frac{(\mchi - \omega)^2}{2M}\Bigr) = \frac{\sigma_F\,\nd_NA^2F_N(q_{\mathrm{abs}})^2}{1 + q_{\mathrm{abs}}/M} ,
\end{equation*}
since the argument of the delta function has slope $1 + (\mchi - \omega)/M$, with $q_{\mathrm{abs}}$ the root of $\mchi = q_{\mathrm{abs}} + q_{\mathrm{abs}}^2/2M$. Expanding, $q_{\mathrm{abs}} = \mchi(1 - \mchi/2M + \dots)$ and $\varGamma_V = \sigma_F\nd_NA^2F_N^2(1 - \mchi/M + \dots)$. Since $\sigma_F \propto \mchi^2$, the rate carries $\mchi^2(1 - \mchi/M)$. The scalar operator weights $\bar S_{nn}$ by $q^2/\mchi^2 = 1 - \mchi/M + \dots$ in addition, and gives $\mchi^2(1 - 2\mchi/M)$: the density row alone, without the longitudinal current that continuity adds, has twice the correction.

In a crystal the free-recoil peak is replaced by the response of \cref{sec:11.3}. Its zero-phonon term carries $\delta(\omega)$ and conserves crystal momentum, $\vect{q} = \vect{G}$. With $\omega = 0$ the absorption requires $|\vect{q}| = \mchi$ for a dark-matter particle at rest, and with speed $v$ the condition $|\vect{G}| = \mchi$ is relaxed only by a relative amount of order $v$: zero-phonon absorption is confined to masses within $\sim10^{-3}$ of the length of a reciprocal-lattice vector. At any other mass only the terms $n \ge 1$ absorb, and in the incoherent approximation \cref{eq:11.3-inc} their weight is $1 - \ee^{-2W_\kappa(q)}$. Along the line $q = \mchi - \omega$ the momentum changes by $\omega$, a fraction $E_R/\mchi = \mchi/2M$ of itself, which rebuilds the factor $1/(1 + q_{\mathrm{abs}}/M)$ in the impulse regime and is negligible below it. Hence \cref{eq:11.4-ferm-nuc}.
